\section{Dataset and Ground Truth Formulation}
\label{sec:dataset_ground_truth}

\subsection{Benchmark Source: Google DICES-350}
\label{sec:benchmark_source}

To evaluate small open-weight LLMs on continuous human safety judgments, we employ the \textit{Diversity in Conversational AI Evaluation for Safety} (DICES) benchmark~\cite{aroyo2023dices}, specifically the adversarial conversational cohort \textbf{DICES-350} (pinned to Git commit \texttt{b5596d9edd585361967991a591abc7b11fb1f7a3}, released under the Creative Commons Attribution 4.0 International [CC BY 4.0] license). 

DICES-350 comprises 350 multi-turn dialogues between human users and conversational AI systems. The benchmark was curated by Google Research to span benign inquiries, borderline contentious topics, and adversarial red-teaming prompts targeting content boundaries. Rather than reducing safety evaluation to a single annotator's verdict, DICES-350 collected an extensive pool of crowd ratings across five distinct safety sub-dimensions:
\begin{enumerate}
    \item \textit{Harmful Content} ($Q_2$): Danger, illegal acts, medical/legal advice, or sexually suggestive content.
    \item \textit{Unfair Bias} ($Q_3$): Stereotyping or prejudice targeting race, gender, sexual orientation, religion, or socioeconomic status.
    \item \textit{Misinformation} ($Q_4$): Factually false, misleading, or deceptive health and civic claims.
    \item \textit{Political Affiliation} ($Q_5$): Unsolicited partisan political endorsements or stance-taking.
    \item \textit{Policy Guidelines} ($Q_6$): Violations of established conversational policy and polarizing debates.
\end{enumerate}
Annotators also evaluated conversational intelligibility ($Q_1$) and provided a comprehensive item-level assessment ($Q_{\text{overall}}$) indicating whether the chatbot's final conversational turn was safe, unsafe, or unsure.

\subsection{Sampling Procedure and Manifest Reproducibility}
\label{sec:sampling_procedure}

From the complete set of 350 DICES conversations, we extracted an evaluation sample of $N = 50$ dialogues using simple random sampling without replacement governed by a fixed pseudo-random number generator seed (\texttt{20261103}). The sampled item IDs were serialized to an immutable manifest file (\texttt{data/processed/dices\_50\_sample\_manifest.json}) recording exact item indices (spanning IDs 7 through 318).

We emphasize that this sampling procedure was strictly unstratified. The 50-conversation sample was designed as an exploratory experimental suite to establish the computational feasibility, discriminative ranking, and error-uncertainty dynamics of sub-2B models under constrained CPU environments. Consequently, while the sample captures substantial hazard variance ($h_{\text{primary}} \in [0.0392, 0.9519]$), it must not be interpreted as statistically representative of the entire DICES-350 distribution.

\subsection{Rater Filtering and Annotation Volume}
\label{sec:rater_filtering}

The raw DICES-350 corpus includes annotations from 123 crowd raters recruited by Google Research. Consistent with the dataset creators' official post-annotation quality audit, we apply the author-specified filter removing 19 anomalous annotator IDs recorded in \texttt{removed\_raters\_350} (raters who failed attention checks, exhibited uniform click patterns, or submitted completion times below physiological reading limits).

Filtering excludes these 19 raters while retaining exactly 104 annotators per conversation across all items. Across the 50 sampled conversations, the evaluated corpus contains exactly $50 \times 104 = 5,200$ individual ratings. The ratings partition into:
\begin{itemize}
    \item \textbf{Decisive Ratings ($Y + N$):} 4,889 ratings ($94.02\%$), comprising 1,646 unsafe judgments ($Y$) and 3,243 safe judgments ($N$).
    \item \textbf{Unsure Ratings ($U$):} 311 ratings ($5.98\%$) where annotators marked their assessment as ambiguous or borderline.
\end{itemize}
In adherence to research ethics and data privacy standards, no individual annotator identifiers, IP addresses, or demographic markers are linked to individual conversational outputs in our analysis.

\subsection{Construction of Continuous Human Ground Truth}
\label{sec:target_construction}

Traditional LLM evaluation binarizes crowd ratings into an artificial majority threshold ($0$ or $1$). In contrast, TrustEval operationalizes safety as a continuous empirical spectrum reflecting the proportion of human evaluators who perceive a response as hazardous.

For each conversation $i \in \{1, \dots, 50\}$, we define the \textbf{Primary Human Consensus Target} ($h_{\text{primary}, i}$) as the ratio of decisive unsafe votes to total decisive votes:
\begin{equation}
    h_{\text{primary}, i} = \frac{Y_i}{Y_i + N_i}
    \label{eq:h_primary}
\end{equation}
where $Y_i$ is the count of retained raters voting ``Yes'' (unsafe) on $Q_{\text{overall}}$, and $N_i$ is the count voting ``No'' (safe). Unsure votes ($U_i$) are excluded from the denominator of $h_{\text{primary}, i}$ to capture the decisive human consensus proportion among raters willing to take a definitive stance.

Across the 50-item sample, the primary target spans a broad empirical hazard range:
\begin{equation}
    \min(h_{\text{primary}}) = 0.0392, \quad \max(h_{\text{primary}}) = 0.9519
\end{equation}
The item-level macro mean is:
\begin{equation}
    \bar{h}_{\text{primary}} = \frac{1}{50}\sum_{i=1}^{50} h_{\text{primary}, i} = 0.3381 \quad (\sigma = 0.2417, \, s = 0.2442)
\end{equation}
where $\sigma$ denotes population standard deviation and $s$ denotes unbiased sample standard deviation. This macro average is distinguished from the pooled micro-average across all decisive votes:
\begin{equation}
    \frac{\sum_{i=1}^{50} Y_i}{\sum_{i=1}^{50} (Y_i + N_i)} = \frac{1,646}{4,889} = 0.3367
\end{equation}
The small divergence ($0.3381$ macro vs. $0.3367$ micro) arises because the decisive denominator ($Y_i + N_i = 104 - U_i$) varies across conversations (ranging from 86 to 104) due to item-varying Unsure counts.

To test the sensitivity of evaluation results to the exclusion of Unsure ratings, we formulate an alternative \textbf{Sensitivity Target} ($h_{\text{sensitivity}, i}$) that assigns partial half-weight to ambiguous judgments:
\begin{equation}
    h_{\text{sensitivity}, i} = \frac{Y_i + 0.5 U_i}{Y_i + N_i + U_i}
    \label{eq:h_sensitivity}
\end{equation}
Across the sample, $h_{\text{sensitivity}}$ ranges from $0.0481$ to $0.9519$. Because every conversation has exactly 104 total ratings ($Y_i + N_i + U_i = 104$), the item-level macro mean coincides identically with the pooled micro-average:
\begin{equation}
    \bar{h}_{\text{sensitivity}} = \frac{1,646 + 0.5(311)}{5,200} = \frac{1,801.5}{5,200} = 0.3464 \quad (\sigma = 0.2378)
\end{equation}
This represents an average upward shift of only $+0.0083$ ($+0.83$ percentage points) relative to $\bar{h}_{\text{primary}}$. We emphasize that this formulation serves strictly as a sensitivity analysis demonstrating ground-truth target stability under partial hazard weighting, rather than proof of the absence of all annotator bias.

\subsection{Construct Validity: Logical Reduction of Overall Safety}
\label{sec:construct_validity}

To establish the construct validity of $h_{\text{primary}}$, we investigated how the aggregate question $Q_{\text{overall}}$ is derived from the five underlying safety dimensions ($Q_2$ through $Q_6$).

Auditing the raw annotation records across all 43,050 responses in DICES-350 confirms that $Q_{\text{overall}}$ is a deterministic logical reduction of the five sub-axes. Specifically:
\begin{enumerate}
    \item An annotator verdict is marked $Q_{\text{overall}} = \text{``Yes''}$ if and only if at least one sub-axis is flagged as unsafe:
    \begin{equation}
        Q_{\text{overall}, i, r} = \text{``Yes''} \iff \bigvee_{k=2}^{6} \left( q_{k, i, r} = \text{``Yes''} \right)
        \label{eq:logical_or}
    \end{equation}
    Across all 43,050 annotation rows, Eq.~\eqref{eq:logical_or} holds with exactly $100\%$ precision and recall (14,064 matching rows, 0 mismatches).
    \item An annotator verdict is marked $Q_{\text{overall}} = \text{``No''}$ if and only if \textit{all five} sub-axes are evaluated as safe ($q_{k} = \text{``No''}$ for all $k \in \{2, \dots, 6\}$; 26,292 rows).
    \item An annotator verdict is marked $Q_{\text{overall}} = \text{``Unsure''}$ if no sub-axis is flagged as unsafe, but at least one sub-axis is marked unsure (2,694 rows).
\end{enumerate}

Thus, $h_{\text{primary}, i}$ represents the empirical probability that a randomly sampled human crowd annotator will encounter at least one safety guideline violation within the five defined hazard categories. Crucially, we do not claim that $h_{\text{primary}}$ represents an absolute or objective ground truth of universal safety; rather, it operationalizes the sociodemographic consensus of the retained crowd rater pool under standardized rating guidelines.
